\section{Two conditional operator estimates}\label{sec:conditional}

The two estimates used in the next section concern an arbitrary
finite-dimensional pure state.
Their constants depend on the dimension of a moved subsystem or of the
support of a local operator, rather than on the dimensions of the remaining
systems.  The first bounds the norm of two overlapping density-matrix
powers.  The second controls how far a local operator expectation departs
from the nonnegative real axis after conjugation by marginal powers.  Later
these estimates
will measure, respectively, the entropy gained and the energy spent when
a small subsystem changes sides in a partition.

For a density matrix $\rho$, write
\[
 \rho^{[z]}=\sum_{p>0}p^z\Pi_p,
 \qquad
 \widehat\rho=\rho+\Pi_{\ker\rho}.
\]
Here $\Pi_p$ is its spectral projection at $p$.  Thus $\rho^{[z]}$ is
zero on the kernel for every complex $z$, including $z=0$, whereas
$\widehat\rho^{iu}$ is unitary.  We omit the brackets on positive real
powers.  On a purification of $\rho$, either convention for an innermost
power gives the same vector.  Tensoring an operator with the identity on
unmentioned factors is understood.
We use the Schatten norms $\|T\|_p=(\Tr|T|^p)^{1/p}$ for
$1\le p<\infty$ and $\|T\|_\infty=\|T\|$; in particular
$\|T\|_2$ is the Hilbert--Schmidt norm.
A density matrix is called faithful if it is positive definite.

\subsection{A relative logarithmic moment bound}

We first record the dimension estimate used in both proofs.  Let $\rho$
be a density matrix on $xU$, let $d=\dim x$, and put
$Q=I_x\otimes\widehat\rho_U$.  If $p_i>0$ and $q_j>0$ are eigenvalues
of $\rho$ and $Q$, respectively, choose eigenbases, with the $Q$ basis
respecting $x\otimes\supp\rho_U$, and give $(i,j)$
probability $p_i|\langle i,j\rangle|^2$.  This defines a real random
variable
\[
 L=\log p_i-\log q_j.
\]
The support inclusion
$\supp\rho\subseteq x\otimes\supp\rho_U$ shows that only eigenvectors
of $Q$ in this latter space contribute.  After restricting $Q$ to
this space and normalizing it by $d$, this is the two-eigenbasis
construction of Nussbaum and
Szko\l a~\cite[Section~3, equation~(11)]{NussbaumSzkola2009}:
the normalized log ratio is $L+\log d$.  We prove the dimension bounds
needed here directly.  They are
\begin{equation}\label{eq:conditional-log-moments}
 \E L=-S(x|U)_\rho,
 \qquad \E e^{L}\le d,
 \qquad \E e^{-L}\le d.
\end{equation}
For completeness, positivity of the blocks of $\rho$ in an $x$-basis
gives, for $v=\sum_{b=1}^d|b\rangle v_b$,
\[
 \langle v,\rho v\rangle
 \le\left(\sum_b\langle v_b,\rho_{bb}v_b\rangle^{1/2}\right)^2
 \le d\sum_b\langle v_b,\rho_Uv_b\rangle.
\]
Thus $\rho\le d(I_x\otimes\rho_U)\le dQ$.  It follows that
$\rho^{1/2}Q^{-1}\rho^{1/2}\le dI$, and hence
$\E e^L=\Tr\rho^2Q^{-1}\le d$.  The other exponential moment equals
$\Tr(\Pi_{\supp\rho}Q)\le d$.  The formula for the mean follows by
taking the two marginal traces.

Put $\ell=\log(ed)$.  Markov's inequality in
\eqref{eq:conditional-log-moments} gives
$\Pr(|L|>v)\le2d e^{-v}$.  Splitting at $2\ell$ and integrating this
tail proves, with universal constants,
\[
 \E\big[L^2e^{|bL|}\big]\le C\ell^2
 \quad\hbox{if } |b|\ell\le c.
\]
Taylor's formula, followed by $\log(1+y)\le y$, therefore yields
\begin{equation}\label{eq:conditional-log-mgf}
 \log\E e^{bL}\le -bS(x|U)_\rho+C b^2\ell^2,
 \qquad |b|\ell\le c.
\end{equation}
These constants are independent of $\dim U$.

A useful consequence is a conditional Schatten estimate.  For
$0<a<1/4$, put
\[
 K=Q^{-a}\rho Q^{-a},\qquad
 r=(1-2a)^{-1},\qquad Z=(\Tr K^r)^{1/r}.
\]
Then, after reducing the universal $c$ if necessary,
\begin{equation}\label{eq:conditional-schatten}
 \log Z\le-2aS(x|U)_\rho+C a^2\ell^2,
 \qquad a\ell\le c.
\end{equation}
Indeed interpolate $\rho^{[z/2]}Q^{-az}$ between operator norm at
$\Re z=0$ and Hilbert--Schmidt norm at $\Re z=r$.  The first bound is
one, and the square of the second is
\[
 \Tr\rho^rQ^{-2ar}=\E e^{(r-1)L}.
\]
The resulting bound at $z=1$ is
\[
 Z=\|\rho^{1/2}Q^{-a}\|_{2r}^2
 \le\big(\E e^{(r-1)L}\big)^{1/r}.
\]
One can obtain this interpolation directly from scalar three-lines:
for a matrix $B=V|B|$ with $\|B\|_{(2r)'}=1$, pair the analytic
matrix with
$|B|^{(2r)'(1-z/(2r))}V^*$ in the trace.  This test has trace norm
one on the first boundary and Hilbert--Schmidt norm one on the second;
at $z=1$ it is $B^*$.  Noninvertible test matrices follow by
approximation.  Now use $r-1=2ar$ in
\eqref{eq:conditional-log-mgf} to prove
\eqref{eq:conditional-schatten}.

\subsection{Moving a subsystem}

\begin{lemma}[Conditional movement estimate]\label{lem:movement}
Let $\theta$ be a unit pure vector on four finite-dimensional systems
$P,x,U,F$, and let $\rho_D$ denote its marginal on $D$.  Set
\[
 d=\dim x,\quad \ell=\log(ed),\quad
 \eta=S(x|P)_\theta+S(x|U)_\theta
      =I(x:F|U)_\theta=I(x:F|P)_\theta.
\]
There are universal constants $c,C>0$ such that, for
$0<a\ell\le c$ and arbitrary density matrices $\sigma$ on $xP$
and $\tau$ on $xU$,
\begin{equation}\label{eq:conditional-movement}
 \left\|\sigma^{a/2}\widehat\rho_P^{-a/2}
          \tau^{a/2}\widehat\rho_U^{-a/2}\theta\right\|
 \le\exp\left(-\frac a2\eta+C a^{5/4}\ell^2\right).
\end{equation}
The same bound holds for positive subnormalized $\sigma,\tau$.
\end{lemma}

\begin{proof}
Purity and strong subadditivity give the displayed identities and
$\eta\ge0$.  It suffices initially to take $\sigma,\tau$ faithful.
Write $d_U=S(x|U)_\theta$ and $d_P=S(x|P)_\theta$.  Consider
\[
 G(z)=\sigma^{az}\widehat\rho_P^{-az}
       \tau^{a(1-z)}\widehat\rho_U^{-a(1-z)}\theta,
 \qquad 0\le\Re z\le1.
\]
The midpoint is the vector in \eqref{eq:conditional-movement}.
All imaginary powers here are unitary.  In particular $G$ is bounded
on this strip, with bounds allowed at this intermediate step to depend
on its fixed invertible factors.

On the left boundary, discard the leftmost unitary factors and move
the remaining imaginary powers past the commuting real powers.  With
$\rho=\rho_{xU}$, $Q=I_x\otimes\widehat\rho_U$, and
$\tau_u=Q^{-iau}\tau Q^{iau}$, this gives
\[
 \|G(iu)\|=M(u):=\|\tau_u^aQ^{-a}\theta\|.
\]
H\"older's inequality and \eqref{eq:conditional-schatten} imply
\[
 M(u)^2=\Tr\tau_u^{2a}K\le Z,
 \qquad \log Z\le-2ad_U+C a^2\ell^2.
\]
We retain the slack in this bound: a large slack already improves
the left boundary, while a small slack will control the distance from
the right-boundary input to $\theta$.  Choose a universal $C_0$ larger than half
the constant in this last bound and set
\begin{equation}\label{eq:conditional-slack}
 \lambda(u)=-ad_U+C_0a^2\ell^2-\log M(u)\ge0,
 \qquad \xi=K^r/\Tr K^r.
\end{equation}
The normalized H\"older overlap then satisfies
\begin{equation}\label{eq:conditional-overlap}
 \Tr\tau_u^{2a}\xi^{1-2a}=M(u)^2/Z\ge e^{-2\lambda(u)}.
\end{equation}

We need to relate this overlap to $\rho$, independently of the
smallest eigenvalues of either marginal.  Taking $\rho$ in place of
$\tau_u$ for this comparison gives
\[
 \|\rho^aQ^{-a}\theta\|
 \ge\langle\theta,\rho^aQ^{-a}\theta\rangle
 =\Tr\rho^{1+a}Q^{-a}
 =\E e^{aL}\ge e^{-ad_U}.
\]
The trace is real and nonnegative.  Thus
$1-\Tr\rho^{2a}\xi^{1-2a}\le C a^2\ell^2$.
The function
$b\mapsto1-\Tr\rho^{[b]}\xi^{[1-b]}$ is nonnegative and concave
on $[0,1]$: expand in their two eigenbases for concavity, and use
H\"older for nonnegativity.  For $2a\le1/2$, concavity and its
nonnegative value at zero give
\[
 1-\Tr\sqrt\rho\sqrt\xi
 \le \frac{1-\Tr\rho^{2a}\xi^{1-2a}}{4a}.
\]
Consequently
\begin{equation}\label{eq:conditional-hellinger}
 H:=\|\sqrt\rho-\sqrt\xi\|_2\le C\sqrt a\,\ell.
\end{equation}

We next show how the slack in \eqref{eq:conditional-slack} controls
the state arriving at the opposite boundary.  Define
\[
 \theta_u=\tau^{-iau}Q^{iau}\theta
          =Q^{iau}\tau_u^{-iau}\theta.
\]
This is a unit vector with the same $P$ marginal as $\theta$.  We
claim
\begin{equation}\label{eq:conditional-state-movement}
 \|\theta_u-\theta\|
 \le C(1+|u|)\big(\sqrt{\lambda(u)}+\sqrt a\,\ell\big).
\end{equation}
Here and in the verification all kernel phases are extended by one.
First,
\[
 \|(Q^{-iau}-\widehat\rho^{-iau})\theta\|^2
 \le a^2u^2\E L^2\le C a^2u^2\ell^2.
\]
For the remaining comparison of $\tau_u$ and $\rho$ insert the phase
of $\xi$.  For any densities $A,B$ and real $b$,
\begin{equation}\label{eq:conditional-phase-hellinger}
 \| (\widehat A^{ib}-\widehat B^{ib})\sqrt B\|_2
 \le C(1+|b|)\|\sqrt A-\sqrt B\|_2.
\end{equation}
To check this, expand in their two spectral bases.  For positive
eigenvalues with ratio $v$, use
$|1-v^{ib}|\le C(1+|b|)|1-\sqrt v|$.  This follows from the
mean-value theorem when $1/4\le v\le4$, and the bound two outside
that interval.  At a zero eigenvalue of $A$ the same assertion
follows from the bound two and $|1-\sqrt0|=1$; zero eigenvalues of
$B$ contribute nothing to the left side.

Replacing $\sqrt\rho$ by $\sqrt\xi$ in the $\tau_u,\xi$ phase
difference costs at most $2H$.  On $\sqrt\xi$, give a pair of
eigenvalues $p$ of $\xi$ and $q$ of $\tau_u$ the probability
$p|\langle p,q\rangle|^2$ and put $v=\log(q/p)$.  Faithfulness of
$\tau_u$ makes this finite, and
\[
 \E e^v\le1,\qquad
 \E e^{2av}\ge e^{-2\lambda(u)}
\]
by \eqref{eq:conditional-overlap}.  On $v\le0$,
\[
 |1-e^{-iauv}|^2\le C(1+u^2)(1-e^{2av}),
\]
and
\[
 \E[(1-e^{2av})1_{v\le0}]
 \le 1-e^{-2\lambda(u)}
       +\E[(e^{2av}-1)1_{v>0}]
 \le 2\lambda(u)+2a.
\]
The last step uses $e^{2av}-1\le2a(e^v-1)$ for $v>0$.
On $v>0$ the squared phase difference is at most
$C a^2u^2 e^v$.  These estimates, followed by
\eqref{eq:conditional-phase-hellinger} for $\xi,\rho$ and
\eqref{eq:conditional-hellinger}, prove
\eqref{eq:conditional-state-movement}.

The phase comparison now lets us recover the other conditional
entropy.  On the right boundary, the remaining $P$ phases merely conjugate
$\sigma$.  Apply \eqref{eq:conditional-schatten} and H\"older on
$xP$ to the input $\theta_u$, whose $P$ marginal is unchanged.
It follows that
\[
 \log\|G(1+iu)\|
 \le-aS(x|P)_{\theta_u}+C a^2\ell^2.
\]
The continuity bound for conditional entropy in
Lemma~\ref{lem:continuity}, using only $\dim x$, and
\eqref{eq:conditional-state-movement} now give
\begin{align*}
 \log\|G(1+iu)\|
 \le{}&-ad_P+C a^2\ell^2\\
 &+C a\ell(1+|u|)^{1/2}
       \big(\lambda(u)^{1/4}+a^{1/4}\ell^{1/2}\big).
\end{align*}
On the left boundary the bound is exactly
$-ad_U+C_0a^2\ell^2-\lambda(u)$.

For clarity, the logarithmic interpolation used here assigns density
\[
 \frac{du}{2\cosh(\pi u)}
\]
to each boundary line when evaluated at the midpoint.  Each line
has total mass $1/2$.  This follows by mapping the strip to the
upper half-plane by $z\mapsto e^{\pi iz}$ and using harmonic
measure at $i$.  Applied to the logarithm of a scalar analytic
function it follows from the submean inequality, first on bounded
domains and then by a limit.  Pairing $G$ with a unit vector in its
midpoint direction gives the vector-norm assertion.  All boundary
norms above are bounded above and away from zero for the fixed
faithful factors, so this application and the limit are legitimate.

Pair the two boundary estimates at the same $u$.  Young's inequality
gives
\[
 -\lambda+C a\ell(1+|u|)^{1/2}\lambda^{1/4}
 \le C\big(a\ell(1+|u|)^{1/2}\big)^{4/3}.
\]
After integration, this term, the quadratic errors, and the remaining
$a^{5/4}\ell^{3/2}$ term are all at most $C a^{5/4}\ell^2$.
The entropy contribution is $-a(d_P+d_U)/2$, proving
\eqref{eq:conditional-movement} for faithful $\sigma,\tau$.
Approximation by faithful densities removes that restriction by
continuity of positive powers.  Finally normalizing subdensities
only increases the left side; the zero subdensity case is immediate.
\end{proof}

\subsection{Comparing marginal phases}

To prove the local conjugation estimate, we first control a change in
products of marginal imaginary powers by conditional mutual information.
For faithful states, Petz's equality theorem characterizes preservation
of relative entropy under restriction to a subalgebra by equality of the
corresponding products of imaginary powers
\cite[Theorem~4(iii),(iv)]{Petz1986}.
The comparison below gives the quantitative estimate needed here, with
constants controlled by $\dim x$.

\begin{lemma}[Marginal-phase comparison]\label{lem:conditional-phases}
Let $\theta$ be a unit pure vector on finite-dimensional systems
$P,x,U,F$, write $\rho_D$ for its marginals, and put
$d=\dim x$ and $\eta=I(x:F|U)_\theta$.  For $\eta>0$ define
\[
 y=\min(1,\eta),\qquad
 \mathcal R_\eta=\sqrt y+\sqrt{\eta\log(ed/y)},
\]
and set $\mathcal R_0=0$.  There is a universal constant $C>0$ such
that, for every $u\in\mathbb R$,
\begin{equation}\label{eq:conditional-modular-bound}
 \big\|\big(\widehat\rho_U^{-iu}\widehat\rho_{xU}^{iu}
       -\widehat\rho_{UF}^{-iu}\widehat\rho_{xUF}^{iu}\big)
                                            \theta\big\|
 \le C|\sinh(\pi u)|\mathcal R_\eta.
\end{equation}
In particular the two phase products agree on $\theta$ if $\eta=0$.
The constants do not depend on $\dim P$, $\dim U$, or $\dim F$.
\end{lemma}

\begin{proof}
The proof writes $\eta$ as the integral of a nonnegative resolvent
defect.  Two moment bounds, involving only $d$, control the ends of
the integral; the entropy controls the part between them.
Initially let
$\rho_f=\rho_{xUF}$ be faithful, and put
\[
 \sigma_f=I_x/d\otimes\rho_{UF},\qquad
 \rho_s=\rho_{xU},\qquad
 \sigma_s=I_x/d\otimes\rho_U,\qquad d=\dim x.
\]
On their Hilbert--Schmidt matrix spaces define positive definite
operators
\[
 A=L_{\sigma_f}R_{\rho_f^{-1}},\qquad
 B=L_{\sigma_s}R_{\rho_s^{-1}},
\]
where $L$ and $R$ denote left and right multiplication.  Define
the isometry
\[
 V(X\sqrt{\rho_s})=(X\otimes I_F)\sqrt{\rho_f}.
\]
It is an isometry by taking the partial trace of $\rho_f$.
Taking matrix elements in the same way with $\sigma_f$ shows
$V^*AV=B$: both sides have the form
$\Tr X^*\sigma_s Y$.  Put $b_0=\sqrt{\rho_s}$ and
$a_0=Vb_0=\sqrt{\rho_f}$.  For $v>0$, let
\[
 \delta_v=(A+vI)^{-1}a_0-V(B+vI)^{-1}b_0.
\]
Expansion of its weighted squared norm, using $V^*(A+vI)V=B+vI$,
gives the resolvent compression identity
\cite[Lemma~2.1]{CarlenVershynina2019}:
\begin{equation}\label{eq:conditional-resolvent-square}
 \|(A+vI)^{1/2}\delta_v\|_2^2
 =\langle a_0,(A+vI)^{-1}a_0\rangle
  -\langle b_0,(B+vI)^{-1}b_0\rangle
 =:h(v)\ge0.
\end{equation}
The logarithmic resolvent formula gives
\begin{equation}\label{eq:conditional-resolvent-entropy}
 \int_0^\infty h(v)\,dv
 =D(\rho_f\|\sigma_f)-D(\rho_s\|\sigma_s)=\eta,
\end{equation}
where $D(\rho\|\sigma)=\Tr\rho(\log\rho-\log\sigma)$.
Indeed $-\langle\sqrt\rho,\log(L_\sigma R_{\rho^{-1}})
\sqrt\rho\rangle=D(\rho\|\sigma)$, and the identity terms in
the log integral cancel.  The last equality in
\eqref{eq:conditional-resolvent-entropy} follows by expanding the
four entropies.

The local domination proved in
\eqref{eq:conditional-log-moments} now reads
$\rho_f\le d^2\sigma_f$.  Therefore
\[
 \langle a_0,A^{-1}a_0\rangle
   =\Tr\rho_f^2\sigma_f^{-1}\le d^2,
 \qquad
 \langle b_0,Bb_0\rangle=\Tr\sigma_s=1.
\]
The first resolvent quadratic form in
\eqref{eq:conditional-resolvent-square} is therefore at most $d^2$.
For another bound use $(A+vI)^{-1}\le v^{-1}I$ and
$(B+vI)^{-1}\ge v^{-1}I-v^{-2}B$.  We obtain
\begin{equation}\label{eq:conditional-resolvent-tail}
 0\le h(v)\le\min(d^2,v^{-2}),\qquad
 \|\delta_v\|_2\le\sqrt{h(v)/v}.
\end{equation}

For $-1<\Re z<0$, scalar integration and spectral calculus give
\[
 A^za_0-VB^zb_0
 =-\frac{\sin\pi z}{\pi}
       \int_0^\infty v^z\delta_v\,dv.
\]
One may approach $z=-iu$ from $-1/4\le\Re z<0$ under this
integral: \eqref{eq:conditional-resolvent-tail} gives an
integrable majorant $d\,v^{-3/4}$ near zero and $v^{-3/2}$ near
infinity.  Hence
\[
 \|A^{-iu}a_0-VB^{-iu}b_0\|_2
 \le C|\sinh\pi u|\int_0^\infty\sqrt{h(v)/v}\,dv.
\]
For $\eta>0$, split the integral at $y/d^2$ and $1/y$.
The two ends cost at most $C\sqrt y$ by
\eqref{eq:conditional-resolvent-tail}; Cauchy--Schwarz and
\eqref{eq:conditional-resolvent-entropy} bound the middle by
\[
 \left(\int h(v)\,dv\right)^{1/2}
 \left(\int_{y/d^2}^{1/y}\frac{dv}{v}\right)^{1/2}
 \le C\sqrt{\eta\log(ed/y)}.
\]
If $\eta=0$, nonnegativity and continuity imply $h(v)=0$ for
all $v>0$, so the discrepancy is zero.  To identify the resulting
vectors, use the definition of $V$:
\[
 \begin{split}
 A^{-iu}a_0
   &=\sigma_f^{-iu}\rho_f^{iu}\sqrt{\rho_f},\\
 VB^{-iu}b_0
   &=\bigl(\sigma_s^{-iu}\rho_s^{iu}\otimes I_F\bigr)
                                      \sqrt{\rho_f}.
 \end{split}
\]
After cancelling the common scalar phase $d^{iu}$ and reversing the
sign of the difference, its Hilbert--Schmidt norm is
\[
 \big\|\big(\widehat\rho_U^{-iu}\widehat\rho_{xU}^{iu}
       -\widehat\rho_{UF}^{-iu}\widehat\rho_{xUF}^{iu}\big)
                                    \sqrt{\rho_f}\big\|_2.
\]
For any operator $T$ on $xUF$ one has
$\|T\theta\|=\|T\sqrt{\rho_f}\|_2$.  This proves
\eqref{eq:conditional-modular-bound} in the faithful case.

For a possibly singular $\rho_f$, replace it by
$(1-\epsilon)\rho_f+\epsilon I/\dim(xUF)$ and use its
reductions in the preceding proof.  Conditional mutual information
converges in finite dimension.  The phase products on the square
root converge in Hilbert--Schmidt norm in the required order.
For example, the innermost lifted phase of $\rho_{xU,\epsilon}$
acts on $\sqrt{\rho_{f,\epsilon}}$.  Its limit lies in
$\supp\rho_{xU}\otimes F$, because
$\supp\rho_f\subseteq\supp\rho_{xU}\otimes F$ and
the square-root mass outside this space tends to zero.  On this
space the inner phase converges, since the mollified marginal
commutes with its limit.  Its output lies in
$x\otimes\supp\rho_U\otimes F$, so the outer
$\rho_{U,\epsilon}$ phase also converges there.  The same argument
applies to the $xUF,UF$ pair.  All phases have norm one, which
controls the vanishing complementary pieces.  Thus the limit is
the product of support phases, equivalently the unitary extensions
appearing in \eqref{eq:conditional-modular-bound}.  No
operator-norm convergence of phases on a limiting kernel is being
used.  The faithful proof was on Hilbert--Schmidt spaces, so it
also applies when a faithful approximation would require a larger
purifying system than $P$.  This establishes
\eqref{eq:conditional-modular-bound} in full generality.
\end{proof}

\subsection{A local conjugation estimate}

We now combine the marginal-phase comparison with an analytic strip
bound for a local positive operator.  This converts the phase control
into the real-power estimate used in the replica argument.

\begin{lemma}[Conditional skew estimate]\label{lem:skew}
Let $\theta$ be a unit pure vector on $P,x,U,F$, put $Y=xU$, and
write $\rho_D$ for its marginals.  Suppose $P=P_0P_1$ and that
$0\le h\le I$ is supported on $P_0x$.  Set
\[
 d_h=(\dim P_0)(\dim x),\qquad \ell_h=\log(ed_h),\qquad
 \eta=I(x:F|U)_\theta=I(x:F|P)_\theta.
\]
Define the entire scalar function
\[
 f(z)=\langle\theta,
       \rho_P^{[z]}\rho_Y^{[-z]}h
       \rho_P^{[-z]}\rho_Y^{[z]}\theta\rangle.
\]
There are universal constants $c,C>0$ such that for
$0<a\ell_h\le c$ and $t=a/2$,
\begin{equation}\label{eq:conditional-skew}
 |f(t)|-\Re f(t)\le C a^2\ell_h^4\eta^{1/8}.
\end{equation}
In particular the left side is zero if $\eta=0$.
\end{lemma}

\begin{proof}
Put $p=\langle\theta,h\theta\rangle\in[0,1]$, so $f(0)=p$.
We first bound $f$ on a strip using the support dimension of $h$,
then compare its imaginary-axis values with $p$, and finally estimate
its real-axis deviation from $[0,\infty)$.  Expand
$h=\sum c\otimes d'$ across $P_0,x$ using matrix units on $P_0$.
There are at most $d_h^2$ terms, with
$\|c\|\|d'\|\le1$ for each.  For $0\le\Re z\le1/4$ the
contribution is the inner product of
\[
 \rho_Y^{[\bar z]}(d')^*\rho_Y^{[-\bar z]}\theta
 \quad\hbox{and}\quad
 \rho_P^{[z]}c\rho_P^{[-z]}\theta.
\]
Each norm is bounded by the norm of its local operator: for a
marginal $\rho$ and $0\le\Re z\le1/2$,
\[
 \|\rho^{[z]}c\rho^{[1/2-z]}\|_2\le\|c\|
\]
by Schatten H\"older.  This is also valid at the endpoints using
support powers.  Reflection
$f(-\bar z)=\overline{f(z)}$ treats negative real parts.  Thus
$|f(z)|\le d_h^2$ on $|\Re z|\le1/4$.  On the imaginary axis
it is an expectation in a unit vector and has absolute value at
most one.  Three-lines on the two half-strips gives
\begin{equation}\label{eq:conditional-narrow-strip}
 |f(z)|\le C\qquad (|\Re z|\le b),
 \qquad b=c_0/\ell_h,
\end{equation}
for a sufficiently small universal $c_0>0$.

On the imaginary axis Schmidt mirroring across $P,YF$ gives
\[
 f(iu)=\langle w_u,hw_u\rangle,
 \qquad
 w_u=\widehat\rho_Y^{iu}\widehat\rho_{YF}^{-iu}\theta.
\]
Compare this with
$w_u^0=\widehat\rho_U^{iu}\widehat\rho_{UF}^{-iu}\theta$.
The latter is a unitary acting only on $UF$, so
$\langle w_u^0,hw_u^0\rangle=p$.  Multiplication of both vectors
first by $\widehat\rho_P^{iu}$ and then by
$\widehat\rho_U^{-iu}$, and a second use of Schmidt mirroring,
show that
\begin{equation}\label{eq:conditional-phase-discrepancy}
 \|w_u-w_u^0\|
 =\| (\widehat\rho_U^{-iu}\widehat\rho_Y^{iu}
       -\widehat\rho_{UF}^{-iu}\widehat\rho_{YF}^{iu})\theta\|.
\end{equation}
Lemma~\ref{lem:conditional-phases} therefore gives
$\|w_u-w_u^0\|\le C|\sinh(\pi u)|\mathcal R_\eta$, where
$\mathcal R_\eta$ is defined there with $d=\dim x$.

It remains to pass from phases to the real-power skew.  Positivity
of $h$ and the identity $\langle w_u^0,hw_u^0\rangle=p$ imply
\[
 |f(iu)-p|\le2\sqrt p\,\|w_u-w_u^0\|
                    +\|w_u-w_u^0\|^2.
\]
In particular, differentiating at zero and using
\eqref{eq:conditional-modular-bound} gives
\begin{equation}\label{eq:conditional-first-derivative}
 f'(0)\in i\R,\qquad |f'(0)|\le C\sqrt p\,\mathcal R_\eta.
\end{equation}
The imaginary character also follows from
$f(-\bar z)=\overline{f(z)}$.

Set $m=\min(1,\mathcal R_\eta)$.  The function
$F(z)=(f(z)-p)e^{z^2}$ has absolute value at most $Cm$ on the
imaginary axis: use \eqref{eq:conditional-modular-bound}, the
Gaussian factor $e^{-u^2}$, and also $|f(iu)-p|\le2$.
On the two outer lines $\Re z=\pm b$ it is bounded by a
universal constant by \eqref{eq:conditional-narrow-strip}.
Three-lines on each half-strip therefore bounds $|F|$ by
$Cm^{1/2}$ on $|\Re z|\le b/2$.  If $m=0$, the identity theorem
already gives $f\equiv p$.  Otherwise Cauchy's estimates on a
circle of radius $b/2$ centered at zero, applied also to the
bounded factor $e^{-z^2}$ on this circle, give
\[
 |f(t)-p-tf'(0)|\le C(t/b)^2m^{1/2},\qquad 0\le t\le b/8.
\]
Write $f'(0)=ib_1$ with $b_1\in\R$.  If $p>0$, then
\[
 |p+itb_1|-p\le t^2b_1^2/p\le C t^2\mathcal R_\eta^2;
\]
if $p=0$, \eqref{eq:conditional-first-derivative} makes $b_1=0$.
Adding the remainder changes $|f|-\Re f$ by at most twice its
absolute value.  Thus
\[
 |f(t)|-\Re f(t)
 \le C a^2\big(\ell_h^2m^{1/2}+\mathcal R_\eta^2\big).
\]
Since $d=\dim x\le d_h$, the quantity supplied by
Lemma~\ref{lem:conditional-phases} satisfies, for $0<\eta\le1$,
$\mathcal R_\eta\le C\sqrt{\eta\log(ed_h/\eta)}$;
for $\eta\ge1$, the bound $\eta\le2\log d\le2\ell_h$ gives
$\mathcal R_\eta\le C\ell_h$.  These two cases imply
\[
 \ell_h^2m^{1/2}+\mathcal R_\eta^2
 \le C\ell_h^4\eta^{1/8},
\]
absorbing the logarithm of $1/\eta$ by the smaller power of
$\eta$.  Taking the universal $c$ in the lemma small enough to
ensure $t=a/2\le b/8$ proves \eqref{eq:conditional-skew}.
\end{proof}

The movement and skew estimates hold without any Hamiltonian
assumption.  A local change of a partition has a nonnegative
conditional information cost, and the departure of a conjugated local
positive expectation from the nonnegative real axis is controlled by
a fractional power of that same cost.  The one-copy constants depend
only on the moved or supporting dimension; the remaining systems may
be arbitrarily large.
